\section{Composición en dos capas: de Additive Heads a Virtual Heads}

El modelo de una capa carece de phase change, mientras que el de dos capas la presenta. La diferencia no está solo en que haya «el doble de heads», sino en que el attention pattern de la segunda capa puede leer las features que escribió la primera, y en que las value paths de ambas capas se combinan en un nuevo effective operator. A continuación se usa path expansion para desarrollar explícitamente esta composición.

\lecturefigure{slide-34-one-layer-formula-recall.jpg}{Repaso de la direct path y las single-head paths del modelo de una capa.}{Video oficial de Stanford Online 00:38:36--00:39:05.}

\subsection{Definición de las matrices OV y QK reutilizables}

Para reducir la notación, el artículo y la clase agrupan las matrices que siempre aparecen juntas:
\[
W_{OV}^{\ell,h}=W_V^{\ell,h}W_O^{\ell,h},
\qquad
W_{QK}^{\ell,h}=W_Q^{\ell,h}(W_K^{\ell,h})^{\top}.
\]
Aquí $\ell$ es el índice de capa y $h$ es el índice de head. La primera describe la escritura de contenido (content write); la segunda, la métrica de compatibilidad (compatibility metric).

\lecturefigure{slide-35-define-wov-wqk.jpg}{Definición de $W_{OV}$ y $W_{QK}$, que expone los dos tipos centrales de circuit de la atención.}{Video oficial de Stanford Online 00:39:06--00:39:33.}

\subsection{Producto de operadores en el modelo de dos capas}

Una vez definidas $W_{QK}$ y $W_{OV}$ para una capa, la pregunta siguiente no es simplemente volver a escribir la segunda capa, sino explicar cómo la segunda capa usa la información que la primera ya escribió en el residual stream. En esta sección se tratan las dos capas como el producto de dos residual-plus-head operators; así se conservan los heads independientes de cada capa y, al mismo tiempo, se hacen visibles las effective paths que solo existen cuando se combinan capas. Este desarrollo es el puente esencial entre «el modelo tiene dos capas» y «el modelo adquiere una nueva capacidad algorítmica».

Si se ignoran LayerNorm/MLP/bias, el operator de cada capa es
\[
\mathcal{L}_{\ell}
=I+\sum_{h\in H_{\ell}}A^{\ell,h}\otimes W_{OV}^{\ell,h}.
\]
La residual transform de dos capas es
\[
X_2=\mathcal{L}_2\mathcal{L}_1X_0.
\]

\lecturefigure{slide-36-two-layer-formula.jpg}{Transformer attention-only de dos capas: producto de dos residual-plus-head operators.}{Video oficial de Stanford Online 00:39:34--00:39:59.}

El Kronecker product cumple la mixed-product identity:
\[
(A_2\otimes W_2)(A_1\otimes W_1)
=(A_2A_1)\otimes(W_2W_1).
\]

\begin{itemize}
  \item $A_2A_1$: composición en dos pasos del position routing;
  \item $W_2W_1$: composición en dos pasos de la content transformation;
  \item elegir la identidad produce paths que omiten alguna capa;
  \item la matriz de atención depende en realidad del residual, por lo que el desarrollo simbólico no elimina la no linealidad del routing.
\end{itemize}

\lecturefigure{slide-37-mixed-product-identity.jpg}{La mixed-product identity separa la composición de posiciones de la composición de contenido.}{Video oficial de Stanford Online 00:40:00--00:40:33.}

\subsection{Los tres tipos de paths tras el desarrollo}

El producto de operadores por sí solo todavía no indica qué rutas de cómputo aparecen realmente. En esta sección se desarrolla $\mathcal{L}_2\mathcal{L}_1$ y se clasifican los términos según omitan alguna capa, pasen por un solo head o atraviesen heads de ambas capas de manera consecutiva. Al leer la expresión siguiente, cada término debe verse como un flujo de información que puede examinarse e intervenirse por separado, en lugar de volver a comprimir todos los summands en una matriz de caja negra; el induction circuit es justamente una instancia concreta del último tipo, la composed path.

El desarrollo da:
\[
I
+\sum_{h\in H_1}A^{1,h}\otimes W_{OV}^{1,h}
+\sum_{k\in H_2}A^{2,k}\otimes W_{OV}^{2,k}
+\sum_{k,h}(A^{2,k}A^{1,h})\otimes
(W_{OV}^{2,k}W_{OV}^{1,h}).
\]

\begin{itemize}
  \item direct path: ambas capas pasan por el residual;
  \item single-head paths: pasan solo por un head de una de las capas;
  \item virtual/composed paths: la salida de un head de la primera capa es leída por un head de la segunda;
  \item un virtual head no agrega parámetros: es el effective circuit de la composición de heads ya existentes.
\end{itemize}

\lecturefigure{slide-38-expanded-two-layer-terms.jpg}{Desarrollo de dos capas: términos direct, de individual head y de composición virtual-head.}{Video oficial de Stanford Online 00:40:34--00:41:39.}

\lecturefigure{slide-39-principled-head-analysis.jpg}{La path expansion ofrece un marco de análisis de attention heads basado en principios.}{Video oficial de Stanford Online 00:41:40--00:42:43.}

\begin{importantbox}{Un circuit no equivale a «un solo head»}
Un algoritmo puede abarcar varios heads y varias capas: un head anterior escribe una feature, un head posterior la lee mediante QK y luego la escribe en los logits mediante OV. Nombrar los mecanismos solo a partir del attention map de un head aislado hace fácil pasar por alto la composición; el path es una unidad más cercana a la ejecución de un programa.
\end{importantbox}

\subsection{Ordenar paths con marginal loss accounting}

En el toy model se puede aplicar ablation a un término de path, o retirarlo, y medir su efecto marginal sobre la loss. La clase presenta los órdenes de magnitud de la direct path, de los heads individuales de la primera y la segunda capa, y de las virtual paths; muchos términos virtuales son pequeños por sí solos, pero no por ello deben descartarse de forma permanente, porque pueden volverse decisivos en modelos más profundos o con ciertos inputs.

\lecturefigure{slide-40-path-loss-contributions.jpg}{Marginal contribution de los distintos path types a la loss.}{Video oficial de Stanford Online 00:42:44--00:42:49.}

\lecturefigure{slide-41-virtual-heads-later.jpg}{En el modelo actual, los términos virtual-head son pequeños, pero pueden ser importantes en modelos más profundos.}{Video oficial de Stanford Online 00:42:50--00:43:05.}

\lecturefigure{slide-42-understand-model-with-twelve-heads.jpg}{Accounting casi completo del toy model: el comportamiento principal se entiende con unos doce heads.}{Video oficial de Stanford Online 00:43:06--00:43:09.}

\lecturefigure{slide-43-head-terms-by-layer.jpg}{Términos de attention heads y reducción de la loss, organizados por layer/type.}{Video oficial de Stanford Online 00:43:10--00:43:21.}

\begin{knowledgebox}{Lectura de la figura: el marginal effect no es el único indicador de importancia}
La head ablation puede quedar compensada por paths redundantes; un head anterior cuyo OV direct logit effect es muy pequeño puede, aun así, alterar con fuerza el QK routing posterior. Conviene examinar a la vez el direct logit effect, la path composition, la dependencia del attention pattern y la ablation específica de cada comportamiento.
\end{knowledgebox}

\subsection{El pattern de la segunda capa depende de lo que escribe la primera}

Aunque la OV path pueda formalizarse como una suma, las $Q$ y $K$ de la segunda capa provienen del residual stream, que contiene los outputs de la primera capa. Por lo tanto, el attention pattern de la segunda capa contiene cross terms: los heads previos cambian «hacia dónde mirar», y no solo «qué escribir después de mirar».

\lecturefigure{slide-44-patterns-read-previous-heads.jpg}{Los attention patterns de la segunda capa usan información de la primera, aunque la OV path simple ignore por el momento la composición.}{Video oficial de Stanford Online 00:43:22--00:43:39.}

\lecturefigure{slide-45-layer-two-qk-context.jpg}{Desarrollo explícito del QK de la segunda capa: tanto queries como keys pueden incluir contribuciones de los heads de la primera capa.}{Video oficial de Stanford Online 00:43:40--00:43:49.}

\lecturefigure{slide-46-positive-eigenvalues-copy.jpg}{Signature OV de copia: los positive eigenvalues hacen más probable que se emita el attended token.}{Video oficial de Stanford Online 00:43:50--00:43:59.}

\lecturefigure{slide-47-second-layer-copying-circuits.jpg}{Varios heads de la segunda capa tienen copying OV circuit, lo que sienta las bases del conditional copy.}{Video oficial de Stanford Online 00:44:00--00:44:11.}

\lecturefigure{slide-48-inspect-patterns-empirically.jpg}{Cuando el QK circuit se vuelve complejo, se infieren mecanismos candidatos a partir de los empirical attention patterns.}{Video oficial de Stanford Online 00:44:12--00:46:07.}

\teachervoice{Olah usa el álgebra, pero también reconoce que un QK circuit complejo se entiende más fácilmente si se empieza por los attention patterns reales. La interpretabilidad mecanicista no admite una sola herramienta: hace que las fórmulas, la visualización, la ablation y los ejemplos se restrinjan entre sí.}

\subsection{Resumen de la sección}

El modelo de dos capas, mediante la multiplicación de operators, produce direct paths, single-head paths y composed virtual paths. Más importante aún, lo que escribe la primera capa modifica el routing de la segunda, de modo que el modelo puede primero desplazar un estado como «el token anterior» y luego, en la capa siguiente, emparejarlo y copiarlo; así obtiene un conditional algorithm del que carece el modelo de una capa.
